% Generated draft: requires mathematical and source review.
\documentclass[11pt]{article}
\usepackage[margin=25mm]{geometry}
\usepackage{fontspec}
\setmainfont{Latin Modern Roman}
\usepackage{amsmath,amssymb,mathtools,unicode-math}
\setmathfont{Latin Modern Math}
\usepackage{enumitem,needspace,xurl,hyperref}
\setlength{\parindent}{0pt}
\setlength{\parskip}{5pt}
\setlength{\emergencystretch}{3em}
\newcommand{\N}{\mathbb{N}}
\newcommand{\Z}{\mathbb{Z}}
\newcommand{\Q}{\mathbb{Q}}
\newcommand{\R}{\mathbb{R}}
\newcommand{\C}{\mathbb{C}}
\newcommand{\F}{\mathbb{F}}
\newcommand{\A}{\mathbb{A}}
\newcommand{\PP}{\mathbb{P}}
\newcommand{\E}{\mathbb{E}}
\newcommand{\eps}{\varepsilon}
\newcommand{\dd}{\,\mathrm{d}}
\begin{document}
\section*{UnsolvedMath ID: 30005639}
\textbf{Source problem code:} catalogue entry
\subsection{Definitions and mathematical statement}
The supplied local registry record identifies this item as UnsolvedMath problem \(30005639\), with source code \texttt{catalogue entry}. Its intended topic is finiteness of a second Galois cohomology group. However, the supplied extraction does not specify that group or its coefficient module, and explicitly marks the statement as \texttt{unrecoverable}. Consequently, no precise finiteness conjecture can be recovered from this record alone.

The excerpt introduces a descending central filtration \(\{F^mG_k\}_{m\geq 1}\), associated to a descending central series \(\{\Pi_{1,1}(m)\}_{m\geq 1}\), and denotes its associated graded quotients by
\[
\operatorname{gr}^mG_k=F^mG_k/F^{m+1}G_k,\qquad m\geq 1.
\]
The underlying field \(k\), the group \(G_k\), the group \(\Pi_{1,1}\), and the construction of this filtration are not supplied. These quotients therefore record the excerpt's notation without specifying a complete arithmetic setting.

\textbf{Unresolved intended statement.} The excerpt says that finiteness of ``the second cohomology group'' is a special case of Jannsen's conjecture, cited there as [Jan89, Conjecture 1]. The supplied statement-verification note further describes this finiteness as a hypothesis in a proved theorem. Neither the theorem nor the relevant cohomology group is included. In schematic notation, a claim of the form
\[
\#H^2(G,M)<\infty
\]
would require identification of \(G\), \(M\), the cohomology theory and topology, and every arithmetic hypothesis and quantifier. This display is only a description of the missing information; it is not a recovered conjecture. In particular, the record does not resolve whether the intended finiteness concerns cardinality or another finiteness property.

\textbf{Additional extracted context.} With omitted preceding notation, the excerpt attributes to [Ish23b, Lemmas 2.11 and 2.12] the assertions that the fixed field
\[
\overline{\Q}^{\ker(\rho_{1,1})}
\]
is a pro-\(p\) extension of \(k(E[p])\), and that
\[
\overline{\Q}^{\ker(\rho_{1,1})}=k(E[p])\,\Omega_k^*,
\]
where the right side denotes a compositum and \(\Omega_k^*\) is described as a pro-\(p\) extension of the mod-\(p\) ray class field \(k(p)\) of \(k\), unramified outside \(p\). The excerpt emphasizes that its second assertion differs from one in a preceding section, and gives as its reason containment of the image of the outer action of \(\operatorname{Aut}(E\setminus O)\) on \(\Pi_{1,1}\) in \(\rho_{1,1}(G_{k(p)})\). The definitions, embeddings, and hypotheses needed to interpret these assertions are absent; they are retained as attributed context, not as independently established statements here.

The excerpt also attributes to Nakamura [Nak95] results interpreting certain Kummer characters from special values of fundamental theta functions as analogues of Soul\'e characters, with certain linear combinations nontrivial. It then displays elliptic Soul\'e characters in the form
\[
\kappa_{m_1,m_2}\colon G_{k(E[p^\infty])}\longrightarrow\Z_p(m_1,m_2).
\]
The extracted expression defining this coefficient twist in terms of \(\chi_1\) and \(\chi_2\) is corrupted, and the index range is missing. The excerpt tentatively suggests applicability of Hain--Matsumoto's technique to prove generation of an unspecified Lie algebra by elements \(\{\sigma_{m_1,m_2}\}\), conditional on nontriviality of the corresponding characters. It also records the speaker's hope of relating the freeness portion of an omitted Conjecture 1 to properties of elliptic multiple zeta values at CM points. Neither passage supplies a complete additional conjecture, and the omitted Conjecture 1 cannot be identified with [Jan89, Conjecture 1] from this record.
\subsection{Short English statement}
UnsolvedMath problem \(30005639\) concerns an unspecified finiteness assertion for second Galois cohomology, attributed to a special case of Jannsen's conjecture. The supplied record marks the formulation as unrecoverable: the cohomology group, coefficients, and hypotheses are missing. Its additional discussion of elliptic Galois representations and Soul\'e characters also depends on omitted definitions. A precise conjecture requires recovery of the surrounding source text and human mathematical review.
\subsection{Sources}
\begin{itemize}
\item Supplied local UnsolvedMath registry record, permanent ID \(30005639\), title ``Jannsen Finiteness for Galois Cohomology,'' source code \texttt{catalogue entry}; collection: Oberwolfach Reports Open Problems. The project identifies \texttt{data/unsolvedmath/problems.json} and the supplementary \texttt{catalog/problems.json} as local registry resources; this draft uses only the supplied complete record.
\item Supplied source citation: \emph{MFO--RIMS Tandem Workshop: Arithmetic Homotopy and Galois Theory (2024)}, Oberwolfach Reports, pp.\ 2377--2488, DOI \texttt{10.4171/owr/2023/42}. Supplied source URL and literature-source link: \url{https://doi.org/10.4171/owr/2023/42}. The record separately gives proposed year 2023; the relationship between that year and the citation's 2024 date is unresolved.
\item Supplied literature source: Jannsen, \emph{l-adic cohomology and Galois cohomology}, \url{https://doi.org/10.1007/978-1-4613-9649-9_5}. The excerpt refers to [Jan89, Conjecture 1], but does not reproduce its statement or full bibliographic entry.
\item Additional internal source labels in the excerpt: [Ish23b, Lemmas 2.11 and 2.12] and Nakamura [Nak95]. Hain--Matsumoto's technique is mentioned without a specific reference. No complete bibliographic entries or links for these references are supplied, and the speaker is not identified.
\item Literature assessment checked at \texttt{2026-08-22}: the supplied curation labels the problem open, identifies the finiteness assertion as a special case of Jannsen's finiteness conjecture, and says only important special cases are known. This is supplied dated curation, not an independently verified literature assessment.
\item The supplied statement-verification note says that the extraction mentions a finiteness hypothesis in a proved theorem without formulating Jannsen's broader conjecture; no separate clean statement is promoted, and the original record remains visible for provenance. The record has statement status \texttt{unrecoverable}, no clean statement, and update timestamp \texttt{2026-09-16T12:06:53Z}. No external source was visited or independently verified for this draft.
\end{itemize}
\end{document}
